(1) 相邻刻线在方向 \(\theta\) 的光程差为 \(d\sin\theta\)，同相主极大满足 \(d\sin\theta_m=m\lambda\)，整数级次还须 \(|m\lambda|\le d\)。真实有效大角度区域可能受刻线包络、偏振及仪器数值孔径限制，因此仅有整数条件并不保证同强度输出。
(2) 对 \(d\sin\theta=m\lambda\) 微分，在 \(m,d\) 固定时得 \(d\cos\theta\,d\theta=m\,d\lambda\)，故 \(d\theta/d\lambda=m/(d\cos\theta)\)。当 \(\cos\theta\to0\) 时几何角分散趋大，但近掠射角通光量、角范围与实际像差将限制可用分辨率。
(3) \(N\) 缝振幅因子为 \(\sin(N\beta)/\sin\beta\)，其中 \(\beta=\pi d\sin\theta/\lambda\)。围绕第 \(m\) 级主峰，相邻第一零点满足相位变量变化 \(|\Delta\beta|=\pi/N\)；换算两近波长的主峰位移刚等于单峰到第一暗纹宽度时，\(\delta\lambda=\lambda/(mN)\)，从而 \(\mathcal R=mN\)，该推导忽略光栅缺陷和像差。
(4) 固定观察角，\(m\lambda=(m+1)\lambda'\)，故相邻较高级次重叠波长 \(\lambda'=m\lambda/(m+1)\)。精确间隔 \(\Delta\lambda=\lambda-\lambda'=\lambda/(m+1)\)，在高阶 \(m\gg1\) 时近似为 \(\lambda/m\)；若改用相邻低一级 \(m-1\)，间隔则是 \(\lambda/(m-1)\)，因此必须说明采用哪种约定。
